\documentclass[11pt]{amsart}
\usepackage[utf8]{inputenc}
\usepackage[margin=1.1in]{geometry}
\usepackage{amsmath,amssymb,amsthm,mathtools}
\usepackage{enumitem}
\usepackage{hyperref}
\hypersetup{colorlinks=true,linkcolor=blue,citecolor=blue,urlcolor=blue,
pdftitle={A comparison inequality for Gaussian vectors and product functions}}

\newtheorem{theorem}{Theorem}[section]
\newtheorem{proposition}[theorem]{Proposition}
\newtheorem{lemma}[theorem]{Lemma}
\newtheorem{corollary}[theorem]{Corollary}
\theoremstyle{definition}
\newtheorem{remark}[theorem]{Remark}

\newcommand{\E}{\mathbb{E}}
\newcommand{\R}{\mathbb{R}}
\newcommand{\Ent}{\operatorname{Ent}}
\newcommand{\cS}{\mathcal{S}}
\newcommand{\Id}{\mathrm{Id}}

\title{A comparison inequality for Gaussian vectors and product functions}
\author{Paata Ivanisvili}
\date{October 28, 2026}

\begin{document}
\sloppy
\begin{abstract}
Let $\Sigma_0$ and $\Sigma_1$ be covariance matrices of Gaussian vectors $\xi=(\xi_1,\dots,\xi_n)$ whose blocks
$\xi_j\in\R^{k_j}$ are standard Gaussian, with $\Sigma_1$ positive definite, and let $p>1$. We prove that
\[
\E_{\Sigma_0}\Big[\prod_jf_j(\xi_j)\Big]\le\Big(\E_{\Sigma_1}\Big[\prod_jf_j(\xi_j)^p\Big]\Big)^{1/p}
\qquad\text{for all measurable }f_j\ge0
\]
whenever some block-diagonal matrix $D$ satisfies $\frac p{p-1}(\Sigma_0-\Sigma_1)\le D\le\Sigma_0$. The necessary
condition is $\Sigma_0\le p\Sigma_1$, and the block-diagonal condition implies it. The two conditions are equivalent
when $\Sigma_0$ lies on the segment between $\Id$ and $\Sigma_1$, and for scalar blocks when, after changing the
signs of some coordinates, all correlations are non-positive and those of $\Sigma_0$ are no larger in absolute value
than those of $\Sigma_1$. For $\Sigma_0=\Id$ the inequality is the reverse H\"older inequality of Chen, Dafnis and
Paouris with equal exponents. In general $\Sigma_0\le p\Sigma_1$ is not sufficient, already for two functions whose
correlation decreases from $\Sigma_0$ to $\Sigma_1$, and Gaussian functions are not extremal: a computer-assisted
example has all Gaussian test functions satisfying the inequality while it fails. Applied to $\Sigma_0=X\Sigma_1X+\Id-X^2$, the results give the forward inequality of
Ivanisvili and Kalantzopoulos with outer exponent $\alpha=1/p$ and unequal parameters $x_j$, a case excluded from
their theorem. Their covariance condition is then necessary and sufficient for two functions with any parameters and
any cross-covariance, and for any number of scalar functions whose correlations become non-positive after sign
changes, provided $x_ix_j\ge0$ whenever the corresponding correlation is non-zero. The proof is a flow along
the segment from $\Sigma_0$ to $\Sigma_1$, with the exponent rising from $1$ to $p$, controlled by a weighted reverse
log-Sobolev inequality. The flow also works along curved paths. In an explicit example, where the block-diagonal
condition holds only for $p\ge3$ and $\Sigma_0\le p\Sigma_1$ holds for $p\ge2.175\ldots$, such a path, with two of the
blocks merged on its first stretch and an exact rational certificate, gives the inequality for $p\ge2.1925$. We show
that these flows stay bounded away from $2.175\ldots$ there; the range in between is open.
\end{abstract}
% --- AI disclosure (site) ---
\noindent\begin{center}\small\itshape
This proof was generated on October 28, 2026, by Grok Bot using Cursor’s agents, guided by several prompts from Paata Ivanisvili
\end{center}\medskip
% --- end AI disclosure ---
\maketitle

\section{Introduction and results}
Let $k_1,\dots,k_n\ge1$ and $K=k_1+\dots+k_n$. We write matrices on $\R^K=\R^{k_1}\times\dots\times\R^{k_n}$ in
blocks $A=(A_{ij})_{i,j=1}^n$, $A_{ij}\in\R^{k_i\times k_j}$, and call $A$ \emph{block-diagonal} if $A_{ij}=0$ for
$i\ne j$. Let $\mathcal C$ be the set of positive semi-definite $K\times K$ matrices whose diagonal blocks are
identities. For $\Sigma\in\mathcal C$, $\E_\Sigma$ denotes expectation with respect to a centred Gaussian vector
$\xi=(\xi_1,\dots,\xi_n)$ with covariance $\Sigma$; so $\xi_j\sim\gamma_{k_j}$, the standard Gaussian measure on
$\R^{k_j}$. For $p>0$ put $\|W\|_{L^p(\Sigma)}=(\E_\Sigma|W|^p)^{1/p}$. Products of values in $[0,\infty]$ are taken
with $0\cdot\infty=0$. The order $A\le B$ is the Loewner order.

\begin{theorem}\label{thm:main}
Let $\Sigma_0,\Sigma_1\in\mathcal C$ with $\Sigma_1$ positive definite, and let $p>1$. Suppose there is a
block-diagonal symmetric matrix $D$ with
\begin{equation}\label{eq:lmi}
\frac p{p-1}\,(\Sigma_0-\Sigma_1)\ \le\ D\ \le\ \Sigma_0 .
\end{equation}
Then for all measurable $f_j:\R^{k_j}\to[0,\infty]$,
\begin{equation}\label{eq:comp}
\E_{\Sigma_0}\Big[\prod_{j=1}^nf_j(\xi_j)\Big]\ \le\ \Big\|\prod_{j=1}^nf_j(\xi_j)\Big\|_{L^p(\Sigma_1)} .
\end{equation}
Conversely, if~\eqref{eq:comp} holds for all $f_j$ of the form $e^{\langle a_j,\cdot\rangle}$, then
$\Sigma_0\le p\Sigma_1$; and~\eqref{eq:lmi} implies $\Sigma_0\le p\Sigma_1$.
\end{theorem}

\begin{remark}\label{rem:basic}
(a) Condition~\eqref{eq:lmi} is a linear matrix inequality in the diagonal blocks of $D$, so its feasibility is a
semidefinite programme. Any such $D$ is positive semi-definite: the diagonal blocks of $\Sigma_0-\Sigma_1$ vanish, so
those of $D$ are $\ge0$.

(b) Let $p^*=\lambda_{\max}(\Sigma_1^{-1/2}\Sigma_0\Sigma_1^{-1/2})$, the smallest $p$ with $\Sigma_0\le p\Sigma_1$.
Since the diagonal entries of $\Sigma_0$ and $\Sigma_1$ agree, $p^*\ge1$, and $p^*=1$ only if $\Sigma_0=\Sigma_1$
(then $\Sigma_1-\Sigma_0\ge0$ has trace $0$). By Lyapunov's inequality, if~\eqref{eq:comp} holds for some $p$, it
holds for every larger $p$. If it holds for every $p>p_0$, it also holds at $p_0$: for $f_j$ with values in
$[\varepsilon,M]$, $0<\varepsilon<M<\infty$, the right side of~\eqref{eq:comp} is continuous in $p$, and general $f_j$
follow by the truncation in Section~\ref{sec:approx}. By Proposition~\ref{prop:holder} below,~\eqref{eq:comp} holds
when $\lambda_{\max}(\Sigma_0)\le p\,\lambda_{\min}(\Sigma_1)$. So the set of $p\ge1$ for which~\eqref{eq:comp} holds is
an interval $[p_c,\infty)$ with $p^*\le p_c\le\lambda_{\max}(\Sigma_0)/\lambda_{\min}(\Sigma_1)$.

(c) Condition~\eqref{eq:lmi} is monotone in $p$: if $D$ satisfies it at $p$ and $p'>p$, then $\frac{c'}cD$ satisfies
it at $p'$, where $c=\frac p{p-1}>c'=\frac{p'}{p'-1}$, because $D\ge0$. The feasible $D$ lie in the compact set
$\{0\le D\le\Sigma_0\}$, so the set of $p>1$ at which~\eqref{eq:lmi} is feasible is closed in $(1,\infty)$. Theorem~\ref{thm:main}
gives $p_c\le p$ for every such $p$, and $p_c=p^*$ when~\eqref{eq:lmi} is feasible at $p=p^*$. The set can be empty
(Remark~\ref{rem:forward}).

(d) If the diagonal blocks $A_j$ of $\Sigma_0$ and $\Sigma_1$ agree but are not identities ($A_j$ positive definite),
the substitution $\xi_j\mapsto A_j^{-1/2}\xi_j$ reduces to $\mathcal C$, since it maps product functions to product
functions. Grouping several blocks into one is allowed when the corresponding parts of $\Sigma_0$ and $\Sigma_1$
agree, and makes~\eqref{eq:lmi} weaker, since $D$ may then have larger blocks.
\end{remark}

Theorem~\ref{thm:path} in Section~\ref{sec:paths} is a version of Theorem~\ref{thm:main} along curved paths, with a
matrix $D$ that may vary along the path; it contains Theorem~\ref{thm:main}. The following elementary bound
complements both. It is a direct consequence of the theorem of
Chen, Dafnis and Paouris~\cite{CDP15}, which we use in the form~\cite[Theorem~B]{IK25}: for $\Sigma\in\mathcal C$ and
$p_j\in\R\setminus\{0\}$, the inequality $\prod_j(\E F_j(\xi_j))^{p_j}\le\E_\Sigma\prod_jF_j(\xi_j)^{p_j}$ holds for all
measurable $F_j:\R^{k_j}\to(0,\infty)$ if and only if $\mathrm{diag}(p_j^{-1}\Id_{k_j})\le\Sigma$, and the reverse
inequality holds for all such $F_j$ if and only if $\mathrm{diag}(p_j^{-1}\Id_{k_j})\ge\Sigma$.

\begin{proposition}\label{prop:holder}
Let $\Sigma_0,\Sigma_1\in\mathcal C$ and $p\ge1$. If there are numbers $s_1,\dots,s_n>0$ with
\begin{equation}\label{eq:holder}
\Sigma_0\ \le\ S\ \le\ p\,\Sigma_1,\qquad S=\mathrm{diag}(s_1\Id_{k_1},\dots,s_n\Id_{k_n}),
\end{equation}
then~\eqref{eq:comp} holds for all measurable $f_j\ge0$. In particular, it holds if
$\lambda_{\max}(\Sigma_0)\le p\,\lambda_{\min}(\Sigma_1)$.
\end{proposition}

\begin{proof}
Let $f_j$ take values in $[\varepsilon,M]$, $0<\varepsilon<M<\infty$, and put $F_j=f_j^{s_j}$. The reverse form of the
theorem quoted above, with $\Sigma=\Sigma_0$ and $p_j=1/s_j$, gives
$\E_{\Sigma_0}\prod_jf_j(\xi_j)\le\prod_j\|f_j\|_{L^{s_j}(\gamma_{k_j})}$, since $S\ge\Sigma_0$. The direct form, with
$\Sigma=\Sigma_1$ and $p_j=p/s_j$, gives $\prod_j\|f_j\|_{L^{s_j}(\gamma_{k_j})}^p\le\E_{\Sigma_1}\prod_jf_j(\xi_j)^p$,
since $\frac1pS\le\Sigma_1$. Together they give~\eqref{eq:comp}; general $f_j\ge0$ follow by the truncation in
Section~\ref{sec:approx}. For the last statement take all $s_j=\lambda_{\max}(\Sigma_0)$: then
$\Sigma_0\le S=\lambda_{\max}(\Sigma_0)\Id\le p\lambda_{\min}(\Sigma_1)\Id\le p\Sigma_1$.
\end{proof}

Proposition~\ref{prop:holder} passes through the independent coupling, so it suits ``forward'' situations, where
correlations decrease from $\Sigma_0$ to $\Sigma_1$; for $\Sigma_1=\Id$ it is sharp (Remark~\ref{rem:forward}).
Theorem~\ref{thm:main} suits ``reverse'' situations, where correlations increase. For instance, for $n=2$, $k_j=1$,
$\mathrm{cov}_{\Sigma_1}(\xi_1,\xi_2)=\rho\in(0,1)$ and $\mathrm{cov}_{\Sigma_0}(\xi_1,\xi_2)=t\rho$ with
$t\in(0,1)$,~\eqref{eq:holder} needs $1+t\rho\le p(1-\rho)$ (by symmetry and convexity one may take $S=s\Id$), while
Corollary~\ref{cor:sharp}(a) below gives~\eqref{eq:comp} under the sharp condition $1-t\rho\le p(1-\rho)$.

\subsection{Cases in which the condition is sharp}
\begin{corollary}\label{cor:sharp}
In each of the following cases,~\eqref{eq:comp} holds if and only if $\Sigma_0\le p\Sigma_1$, for every $p>1$.
\begin{enumerate}[label=(\alph*)]
\item $\Sigma_1\in\mathcal C$ is positive definite and $\Sigma_0=(1-t)\Id+t\Sigma_1$ with $t\in[0,1]$. Here
$\Sigma_0\le p\Sigma_1$ reads $1-t+t\lambda\le p\lambda$, $\lambda=\lambda_{\min}(\Sigma_1)$.
\item All $k_j=1$, $\Sigma_1$ is positive definite, and there are signs $\varepsilon_j\in\{\pm1\}$ with
\begin{equation}\label{eq:signs}
\varepsilon_i\varepsilon_j(\Sigma_1)_{ij}\ \le\ \varepsilon_i\varepsilon_j(\Sigma_0)_{ij}\ \le\ 0\qquad(i\ne j).
\end{equation}
\item $n=2$, $\Sigma_1=\begin{psmallmatrix}\Id&C\\C^T&\Id\end{psmallmatrix}$ is positive definite and
$\Sigma_0=\begin{psmallmatrix}\Id&tC\\tC^T&\Id\end{psmallmatrix}$ with $t\le0$ and $|t|\,\|C\|\le1$. Here
$\Sigma_0\le p\Sigma_1$ reads $(p-t)\|C\|\le p-1$.
\end{enumerate}
At $p=1$ both~\eqref{eq:comp} and $\Sigma_0\le\Sigma_1$ mean $\Sigma_0=\Sigma_1$: exponential $f_j$ in~\eqref{eq:comp}
give $\Sigma_0\le\Sigma_1$, and then the trace argument of Remark~\ref{rem:basic}(b) applies.
\end{corollary}

For $t=0$, (a) is the reverse H\"older inequality of Chen, Dafnis and Paouris~\cite{CDP15} with equal exponents
(\cite[Theorem~B]{IK25} with $p_j=p$): $\prod_j\E f_j\le\|\prod_jf_j(\xi_j)\|_{L^p(\Sigma_1)}$ iff
$p\lambda_{\min}(\Sigma_1)\ge1$. Inequalities of this direction are also called inverse Brascamp--Lieb
inequalities~\cite{BW22}; Barthe's reverse Brascamp--Lieb inequality~\cite{Bar98} can be obtained from them by a
limiting argument~\cite{IK25}. For $t=1$, (a) is Lyapunov's inequality. Case (a) contains every pair with $n=2$ of
the form $\Sigma_1=\begin{psmallmatrix}\Id&C\\C^T&\Id\end{psmallmatrix}$,
$\Sigma_0=\begin{psmallmatrix}\Id&tC\\tC^T&\Id\end{psmallmatrix}$, $t\in[0,1]$, since then $\Sigma_0=(1-t)\Id+t\Sigma_1$;
together with (c), every $t\le1$ is covered. For some $t>1$ the condition $\Sigma_0\le p\Sigma_1$ is not sufficient
(Proposition~\ref{prop:notsuff}). Condition~\eqref{eq:signs} says that after changing the signs of
some coordinates all correlations are non-positive, and those of $\Sigma_0$ are no larger in absolute value than
those of $\Sigma_1$. Signs making the off-diagonal entries of $\Sigma_1$ non-positive exist if and only if every
cycle in the graph of its non-zero off-diagonal entries contains an even number of positive entries. (Along a cycle
the product of the signs $-\varepsilon_i\varepsilon_j$ is $(-1)^{\text{length}}$; conversely, choose the signs along a
spanning tree of each component.) This holds for instance when $n=2$, when the graph is a forest, and when all
correlations are non-positive. Three variables with pairwise positive correlations are a non-example.

By Remark~\ref{rem:basic}(d), (a) and (c) also apply to grouped blocks. Let $\Sigma_1$ be positive definite, and
suppose the index set is partitioned into groups such that $\Sigma_0$ and $\Sigma_1$ agree on each group, and $(\Sigma_0)_{ij}=t(\Sigma_1)_{ij}$ whenever $i$ and
$j$ lie in different groups. If $t\in[0,1]$, or if $t<0$, there are two groups and $\Sigma_0\ge0$, then~\eqref{eq:comp} holds
if and only if $\Sigma_0\le p\Sigma_1$. Indeed, after the normalisation of Remark~\ref{rem:basic}(d) on each group
(a congruence, which preserves the Loewner order), $\Sigma_0$ becomes $(1-t)\Id+t\Sigma_1$, respectively the matrix of
(c).

\subsection{The forward inequality with outer exponent \texorpdfstring{$1/p$}{1/p}}
For $x\in[-1,1]$ let $T_xf(y)=\E f(xy+\sqrt{1-x^2}\,G)$, $G\sim\gamma_k$, be the Ornstein--Uhlenbeck (Mehler)
operator. Ivanisvili and Kalantzopoulos~\cite[Thm.~9]{IK25} consider, for $\Sigma\in\mathcal C$, $r_j\in[-1,1]$, $p_j\ne0$
and $\alpha\in\R\setminus\{0\}$, the forward inequality
\begin{equation}\label{eq:IK}
\Big\|\prod_j\big(T_{r_j}f_j(\xi_j)\big)^{p_j}\Big\|_{L^\alpha(\Sigma)}\le\Big\|\prod_jf_j(\xi_j)^{p_j}\Big\|_{L^1(\Sigma)}
\end{equation}
for positive $f_j$. They show that for $\alpha\in(-\infty,0)\cup[1,\infty)$ it holds iff
\begin{equation}\label{eq:IKcond}
\mathrm{diag}\Big(\frac{1-r_j^2}{p_j}\Id_{k_j}\Big)\le\big((1-\alpha r_ur_v)\Sigma_{uv}\big)_{u,v=1}^n
\end{equation}
(\cite[(2.16), (2.17)]{IK25}). The range $\alpha\in(0,1)$ is excluded. The companion note~\cite{FR} settles it for
equal parameters $r_j=x$ and equal exponents $p_j=p$ (\cite[Thm.~1.1]{FR} for $p\ge1$, and~\cite[Cor.~1.4]{FR} for all
$p\ne0$ when $\Sigma$ is non-singular), and for two $\rho$-correlated functions when $\alpha=1/p$, for all
$r_1,r_2$~\cite[Prop.~6.1]{FR}. It also shows that for two functions with unequal parameters and any $\alpha p>2$,
\eqref{eq:IKcond} is not sufficient in general~\cite[Props.~6.3, 6.5]{FR}, while for $\alpha p>1$ and positive
$\rho,r_1,r_2$ a flow with separate noise parameters gives a sufficient condition under which it is
sharp~\cite[Thm.~6.9]{FR}.

\begin{corollary}\label{cor:IK}
Let $\Sigma\in\mathcal C$ be positive definite, $p>1$, $x_j\in[-1,1]$, and $X=\mathrm{diag}(x_j\Id_{k_j})$. Then
\eqref{eq:IK} with $r_j=x_j$, $p_j=p$ and $\alpha=1/p$ is the inequality~\eqref{eq:comp} for
$\Sigma_0=X\Sigma X+\Id-X^2$ and $\Sigma_1=\Sigma$, and~\eqref{eq:IKcond} is $\Sigma_0\le p\Sigma$. So
\eqref{eq:IKcond} is necessary, and it is sufficient whenever~\eqref{eq:lmi} or~\eqref{eq:holder} is feasible. In
particular, \eqref{eq:IK}$\Leftrightarrow$\eqref{eq:IKcond} in each of the following cases:
\begin{enumerate}[label=(\roman*)]
\item $x_ix_j=t\ge0$ for all $i\ne j$ with $\Sigma_{ij}\ne0$; for instance, all $x_j$ are equal;
\item all $k_j=1$, $x_ix_j\ge0$ whenever $\Sigma_{ij}\ne0$, and some signs $\varepsilon_j$ satisfy
$\varepsilon_i\varepsilon_j\Sigma_{ij}\le0$ for all $i\ne j$;
\item $n=2$, with any $x_1,x_2$, any $k_1,k_2$ and any cross-covariance $\Sigma_{12}$;
\item $x_j=1$ for all $j$ in a set $G$ of indices, and either $x_j=0$ for all $j\notin G$, or $G$ contains all
indices but one.
\end{enumerate}
\end{corollary}

In case (i), $\Sigma_0=(1-t)\Id+t\Sigma$ is the noised covariance for the equal parameter $x=\sqrt t$, so (i) is
the case $r=1$ of~\cite[Thm.~1.1]{FR}, which allows arbitrary block sizes. Case (iii) with $x_1x_2\ge0$ is contained in (i), and
case (iv) with a non-negative parameter ($x_{j_0}$, or $0$) follows from the same theorem after grouping and
normalising as in Remark~\ref{rem:basic}(d); with a negative parameter both follow from~\cite{CDP15} through
Proposition~\ref{prop:holder}. For two $\rho$-correlated functions case (iii) is~\cite[Prop.~6.1]{FR}. We include (iii) and (iv) for
completeness. Case (ii) with $n\ge3$ and genuinely unequal parameters is new as far as we know.

\subsection{Limits of the method}
\begin{proposition}\label{prop:example}
Let $n=3$, $k_j=1$, $\Sigma_1$ the equicorrelated matrix with correlation $\frac12$, and $x=(1,\frac12,0)$, so
$\Sigma_0=X\Sigma_1X+\Id-X^2=\begin{psmallmatrix}1&1/4&0\\1/4&1&0\\0&0&1\end{psmallmatrix}$. Then
$p^*=\frac{11+\sqrt{41}}8=2.175\ldots$. Condition~\eqref{eq:lmi} is feasible if and only if $p\ge3$, and
condition~\eqref{eq:holder} if and only if $p\ge\frac52$. So~\eqref{eq:comp} holds for $p\ge\frac52$, while the
necessary condition only requires $p\ge p^*$.
\end{proposition}

Proposition~\ref{prop:certificate} below improves $\frac52$ to $2.1925$, and we do not know whether~\eqref{eq:comp}
holds for $p\in[p^*,2.1925)$ here. The example is not covered by
Corollaries~\ref{cor:sharp} and~\ref{cor:IK}: all correlations of $\Sigma_1$ are positive, so~\eqref{eq:signs} fails,
and $\Sigma_0$ and $\Sigma_1$ agree on no group of two or more coordinates. Numerically, no Gaussian test
function violates~\eqref{eq:comp} at $p=p^*$ (Section~\ref{sec:num}), and by the following remark~\eqref{eq:comp} holds
near constant functions, in the uniform norm, for every $p>p^*$.

\begin{remark}\label{rem:second}
Let all $k_j=1$, let $(h_d)_{d\ge0}$ be the normalised Hermite polynomials, and let $g_j$ be bounded with
$\E g_j(\xi_j)=0$, so that $g_j=\sum_{d\ge1}a_{j,d}h_d$ in $L^2(\gamma_1)$. Put $a_d=(a_{1,d},\dots,a_{n,d})$ and let
$\Sigma^{\circ d}$ be the entrywise $d$-th power. For $f_j=1+\varepsilon g_j$, the right side of~\eqref{eq:comp} minus
the left side equals
\[
\frac{\varepsilon^2}2\sum_{d\ge1}a_d^T\big(p\Sigma_1^{\circ d}-\Sigma_0^{\circ d}\big)a_d+O(\varepsilon^3)
\qquad(\varepsilon\to0).
\]
Indeed, $\E_{\Sigma_0}\prod_j(1+\varepsilon g_j)=1+\varepsilon^2\sum_{i<j}\E_{\Sigma_0}[g_ig_j]+O(\varepsilon^3)$, and
$N(\varepsilon)=\E_{\Sigma_1}\prod_j(1+\varepsilon g_j)^p$ is smooth near $0$ with $N(0)=1$, $N'(0)=0$ and
$N''(0)=p(p-1)\sum_j\E g_j^2+2p^2\sum_{i<j}\E_{\Sigma_1}[g_ig_j]$, so that
$N^{1/p}=1+\frac{\varepsilon^2}{2p}N''(0)+O(\varepsilon^3)$. Mehler's formula gives
$\E_\Sigma[g_i(\xi_i)g_j(\xi_j)]=\sum_da_{i,d}a_{j,d}\Sigma_{ij}^d$. In particular $p\Sigma_1^{\circ d}\ge\Sigma_0^{\circ d}$
for all $d\ge1$ is necessary for~\eqref{eq:comp} (truncate $h_d$ as in~\cite[Prop.~6.2]{FR}); $d=1$ is
$\Sigma_0\le p\Sigma_1$. In Proposition~\ref{prop:example} the off-diagonal entries of $p\Sigma_1^{\circ d}-\Sigma_0^{\circ d}$
are $p2^{-d}-4^{-d}$, $p2^{-d}$ and $p2^{-d}$, so for $d\ge2$ and $p>2$ Gershgorin's theorem gives
$\lambda_{\min}(p\Sigma_1^{\circ d}-\Sigma_0^{\circ d})\ge p-1-2p\,2^{-d}\ge\frac p2-1>0$. For $p>p^*$ also
$p\Sigma_1-\Sigma_0>0$. Hence for every $p>p^*$ the quadratic form is at least $c\sum_j\E g_j^2$ with $c>0$. The remainder is uniform:
if $\|g_j\|_\infty\le M$ and $|\varepsilon|M\le\frac12$, it is at most $C|\varepsilon|^3M\sum_j\E g_j^2$ with
$C=C(n,p)$, because every term of order $k\ge3$ contains two factors $g_i,g_j$ whose product is bounded by
$\frac12(g_i^2+g_j^2)$ and $k-2$ factors bounded by $M$. So for every $p>p^*$ there is $\delta(p)>0$ such
that~\eqref{eq:comp} holds whenever $\|f_j/\E f_j-1\|_\infty\le\delta(p)$ for all $j$: a violation for some
$p\in(p^*,2.1925)$, if there is one, must be far from constants in the uniform norm.
\end{remark}

\begin{proposition}\label{prop:certificate}
In the example of Proposition~\ref{prop:example},~\eqref{eq:comp} holds for every $p\ge\frac{877}{400}=2.1925$.
\end{proposition}

\begin{proof}[Proof (computer-assisted, exact arithmetic)]
Let $b=\frac9{40}$, $C_1=\begin{psmallmatrix}1&1/4&b\\1/4&1&b\\b&b&1\end{psmallmatrix}$ and $p_1=\frac{139783}{100000}$.
The matrices $\Sigma_0$ and $C_1$ agree on the groups $\{1,2\}$ and $\{3\}$, and $\Sigma_0$ vanishes between them, so
by the grouping paragraph after Corollary~\ref{cor:sharp} (case (a) with $t=0$),~\eqref{eq:comp} holds from $\Sigma_0$
to $C_1$ with exponent $p_1$ as soon as $\Sigma_0\le p_1C_1$; this is checked exactly (the threshold is
$1/(1-\sqrt{8/5}\,b)=1.3978291\ldots$). Next, Theorem~\ref{thm:path} applies to a piecewise-linear path through
rational correlation matrices $K_0=C_1,K_1,\dots,K_{256}=\Sigma_1$, with scalar blocks, pieces of length
$\ell=\frac{4501}{2560000}$ and a rational diagonal matrix $D_k$ on the $k$-th piece. On a piece $K(s)-D_k$ is affine in
$s$ and $K'(s)=(K_{k+1}-K_k)/\ell$, so~\eqref{eq:pathcond} holds there if $D_k\ge0$, $K_k-D_k\ge0$, $K_{k+1}-D_k\ge0$ and
$(K_{k+1}-K_k)+\ell D_k\ge0$; positive definiteness of all $K_k$ gives that of $K(s)$. These $1281$ conditions are
verified exactly (principal minors) by \texttt{numerics/comparison\_chain\_certificate.py}, which stores the rational
data and re-checks them without any solver. This gives~\eqref{eq:comp} from $C_1$ to $\Sigma_1$ with exponent
$e^{256\ell}=e^{0.4501}$. Applying it to the functions $f_j^{p_1}$ after the first step, the exponents multiply:
\[
\E_{\Sigma_0}\Big[\prod_jf_j\Big]\le\Big(\E_{C_1}\Big[\prod_jf_j^{p_1}\Big]\Big)^{1/p_1}
\le\Big(\E_{\Sigma_1}\Big[\prod_jf_j^{p_1e^{0.4501}}\Big]\Big)^{1/(p_1e^{0.4501})} .
\]
Finally $p_1e^{0.4501}=2.19245\ldots<\frac{877}{400}$ in ball arithmetic, and Lyapunov's inequality gives every
$p\ge\frac{877}{400}$.
\end{proof}

The first step is the segment flow of Theorem~\ref{thm:main} for the groups, with $D=\Sigma_0$ (for this $D$,
condition~\eqref{eq:lmi} is exactly $\Sigma_0\le p_1C_1$), so by Remark~\ref{rem:path}(a) the whole construction is one
admissible path for Theorem~\ref{thm:path} with grouped blocks on its first piece. The data were found by
semidefinite programming (Section~\ref{sec:num}); this construction was suggested by an independent referee of this
note. Optimising $b$ gives $2.1949$ and $2.1931$ with $64$ and $128$ scalar pieces, which suggests about $2.191$ in the
limit. For comparison (all numerical): along the straight segment, letting $\tilde D$ vary as in
Remark~\ref{rem:path}(d) already gives $2.3148$ instead of $3$
(\texttt{numerics/comparison\_segment\_varying\_D.py}); curved paths with scalar blocks give at best about
$2.287$ ($2.2957$, $2.2913$ and $2.2891$ for $64$, $128$ and $256$ pieces; $\frac{229}{100}$ is certified the same way by
\texttt{numerics/comparison\_path\_certificate.py 256 1657/2000 229/100}), since the optimal scalar paths are not
straight: the correlation of $\xi_1$ and $\xi_2$ first decreases slightly. A discretised first stretch with frozen
$c_{12}$ gives $\frac{221}{100}$ (\texttt{numerics/comparison\_multistage\_certificate.py GS 32 0.4,0.6 7929/10000
221/100}). So most of the gain over Theorem~\ref{thm:main} comes from letting $\tilde D$ vary, and the rest from merging
blocks. By Proposition~\ref{prop:away}, no such flow reaches $p^*$ here, nor any exponent close to it; numerically they
bottom out near $2.191$.

The second variation also shows that the necessary condition is not sufficient in general.

\begin{proposition}\label{prop:notsuff}
Let $n=2$, $k_j=1$, $\mathrm{cov}_{\Sigma_1}(\xi_1,\xi_2)=\rho\in[0,1)$ and
$\mathrm{cov}_{\Sigma_0}(\xi_1,\xi_2)=\theta\in(\rho,1)$. Then $p^*=\frac{1+\theta}{1+\rho}$, and~\eqref{eq:comp} fails for
every $p<\sup_{d\ge1}\frac{1+\theta^d}{1+\rho^d}$. If $(1+\theta)(1+\rho)>2$, this supremum exceeds $p^*$ (already the
term $d=2$ does), so $\Sigma_0\le p\Sigma_1$ does not imply~\eqref{eq:comp}.
\end{proposition}

\begin{proof}
$p\Sigma_1-\Sigma_0\ge0$ iff $|p\rho-\theta|\le p-1$. If $p\rho\le\theta$ this reads $p\ge\frac{1+\theta}{1+\rho}$; if
$p\rho>\theta$ it reads $p(1-\rho)\ge1-\theta$, which holds for all $p\ge1$. Since $\frac{1+\theta}{1+\rho}\le\frac\theta\rho$,
$p^*=\frac{1+\theta}{1+\rho}$. By Remark~\ref{rem:second},~\eqref{eq:comp} requires $p\Sigma_1^{\circ d}\ge\Sigma_0^{\circ d}$
for every $d$, and the vector $(1,1)$ gives $p-1+p\rho^d-\theta^d\ge0$, i.e.\ $p\ge\frac{1+\theta^d}{1+\rho^d}$. Finally
$(1+\theta^2)(1+\rho)-(1+\theta)(1+\rho^2)=(\theta-\rho)\big((1+\theta)(1+\rho)-2\big)$, which is positive under the
hypothesis.
\end{proof}

For $\rho=0$ the hypothesis cannot hold, in accordance with the forward H\"older inequality. For instance, for
$\theta=0.9$ and $\rho=0.5$ the supremum is attained at $d=4$, and~\eqref{eq:comp} fails on
$[\frac{19}{15},1.5587\ldots)$. By Remark~\ref{rem:second}, $f_1=f_2=1+\varepsilon(y^2-1)/\sqrt2$ violates it for every
$p<1.448$ once $\varepsilon$ is small enough (depending on $p$); the expansion there also holds for this unbounded
perturbation, since $f_j\ge1-\varepsilon/\sqrt2>0$ and all moments are finite; quadrature confirms this at $p=\frac{19}{15}$ and
$p=1.35$ (\texttt{numerics/comparison\_hermite\_scan.py}). Gaussian test functions do not detect the
higher levels: for these $\theta,\rho$, a numerical optimisation over all $f_j=e^{a_jy-b_jy^2/2}$ finds no violation
for $p\ge1.497$. The following proposition makes this rigorous at another point. The Hermite levels do not
determine the threshold either: at $p=1.56$, where all levels are positive semi-definite (smallest margin
$0.0014$, at $d=4$), an optimised non-Gaussian pair $f_1,f_2$ violates~\eqref{eq:comp} by $4\cdot10^{-6}$ in the
logarithm, stably under refinement of the quadrature, while nothing is found for $p\ge1.6$
(\texttt{numerics/comparison\_n2\_forward\_attack.py}; floating point). So for $\theta=0.9$, $\rho=0.5$
there is a violation at $p=1.56$ and none was found for $p\ge1.6$ (floating point), while Proposition~\ref{prop:holder}
only gives $p\ge3.8$ and Theorem~\ref{thm:path} $p\ge5$.

In the setting of Corollary~\ref{cor:IK} the correlations of $\Sigma_0=X\Sigma X+\Id-X^2$ are no larger in absolute
value than those of $\Sigma$, so Proposition~\ref{prop:notsuff} does not apply there. In $196515$ random instances of
that setting with $n\in\{3,4\}$ all levels $2\le d\le8$ of Remark~\ref{rem:second} were positive semi-definite at $p^*$
(\texttt{numerics/comparison\_hermite\_scan.py}), and a targeted minimisation of the smallest level margin for
$n\le5$, with $\lambda_{\min}(\Sigma)\ge0.05$, found only positive values
(\texttt{numerics/comparison\_hermite\_target.py}).

\begin{proposition}\label{prop:gauss}
Let $n=2$, $k_j=1$, $\mathrm{cov}_{\Sigma_1}(\xi_1,\xi_2)=\frac35$, $\mathrm{cov}_{\Sigma_0}(\xi_1,\xi_2)=\frac{93}{100}$ and
$p=\frac32$. Then every Gaussian test function $f_j(y)=c_je^{a_jy-b_jy^2/2}$ ($c_j>0$, $a_j,b_j\in\R$)
satisfies~\eqref{eq:comp}, but~\eqref{eq:comp} fails for some measurable $f_j$ with values in $[\frac12,\frac32]$. So
Gaussian functions are not extremal for~\eqref{eq:comp}.
\end{proposition}

Inequalities of Brascamp--Lieb type bound a multilinear Gaussian form by a product of norms of the individual
functions, and under the hypotheses of~\cite{Lieb90,BW22,CL21,NT25} Gaussian functions are extremal for them.
Inequality~\eqref{eq:comp} compares two multilinear forms; unless $\Sigma_1=\Id$ its right side is not a product of
norms, and these results do not apply. The functions in Proposition~\ref{prop:gauss} are not log-concave; whether
Gaussian functions are extremal for~\eqref{eq:comp} among log-concave $f_j$, the setting of~\cite{MNT26}, is open.

\begin{proof}[Proof (computer-assisted: exact arithmetic and ball arithmetic)]
Write $\theta=\frac{93}{100}$, $\rho=\frac35$. Here $p^*=\frac{193}{160}<\frac32$, and
$(p-1)+p\rho^5-\theta^5=-\frac{790483693}{10^{10}}<0$, so~\eqref{eq:comp} fails by Proposition~\ref{prop:notsuff} (level
$d=5$; the bounded perturbations $1+\varepsilon g$ of Remark~\ref{rem:second} take values in $[\frac12,\frac32]$ for
small $\varepsilon$).

Let $B=\mathrm{diag}(b_1,b_2)$, $K_0=\Sigma_0^{-1}+B$ and $K_1=\Sigma_1^{-1}+pB$. If $K_1$ is not positive definite,
the right side of~\eqref{eq:comp} is infinite. Otherwise $K_0\ge\frac1pK_1>0$, because $\Sigma_0\le p\Sigma_1$, and
Gaussian integrals give
\[
\log\frac{\E_{\Sigma_0}[f_1f_2]}{\|f_1f_2\|_{L^p(\Sigma_1)}}=\frac12\Big(\frac1p\log\det(\Id+p\Sigma_1B)
-\log\det(\Id+\Sigma_0B)\Big)+\frac12a^T\big(K_0^{-1}-pK_1^{-1}\big)a ,
\]
where the last term is $\le0$. With $A(b)=\det(\Id+\Sigma_0B)=1+b_1+b_2+\frac{1351}{10^4}b_1b_2$ and
$\tilde B(b)=\det(\Id+p\Sigma_1B)=1+\frac32(b_1+b_2)+\frac{36}{25}b_1b_2$, it remains to show
$P=A^3-\tilde B^2\ge0$ on $\mathcal D=\{b:K_1>0\}$. In the variables $u_i=b_i+\frac{25}{24}$,
$\mathcal D=\{u_1,u_2>0,\ u_1u_2>\frac{25}{64}\}$, $\tilde B=\frac{36}{25}u_1u_2-\frac9{16}$, and
$A=\kappa u_1u_2+\frac{8249}{9600}(u_1+u_2)-\frac{8633}{9216}$ with $\kappa=\frac{1351}{10^4}$. Since
$u_1+u_2\ge\frac54$ on $\mathcal D$, $A\ge\kappa u_1u_2+0.137$ there. On the boundary $\tilde B=0$, so $P=A^3>0$.
Along any sequence in $\mathcal D$ with $|u|\to\infty$, $P\to\infty$: where $u_1u_2$ stays bounded, $\tilde B$ stays bounded
and $A\to\infty$; where $u_1u_2\to\infty$, $P\ge\kappa^3(u_1u_2)^3-(\frac{36}{25}u_1u_2)^2$. So if $\inf_{\mathcal D}P\le0$,
it is attained at a critical point of $P$ in $\mathcal D$. The partial derivatives of $P$ are coprime, and a
lexicographic Gr\"obner basis of the ideal they generate is $\{b_2+d(b_1),\,h(b_1)\}$ with $\deg h=9$. The
polynomial $h$ has exactly three real roots, isolated in exact arithmetic. They give the critical points $(0,0)$,
where $P=0$; $b\approx(-0.3024,-0.3024)\in\mathcal D$, where $P\in0.0173028661872\pm2\cdot10^{-14}$; and
$b\approx(-39.83,-39.83)\notin\mathcal D$. Hence $P\ge0$ on $\mathcal D$
(\texttt{numerics/certify\_gaussian\_not\_extremal.py}).
\end{proof}

\section{Preliminaries}
Let $\cS$ be the set of $C^\infty$ functions $f$ on some $\R^k$ with $\inf f>0$ such that $f$ and all its derivatives
are bounded. If $g$ is measurable with values in $[\varepsilon,M]$, $0<\varepsilon<M<\infty$, and $s\in[0,1)$, then
$T_sg\in\cS$: differentiate the Gaussian density in $T_sg(y)=\int g(w)\phi_{1-s^2}(w-sy)\,dw$ for $s>0$, where
$\phi_v$ is the density of $N(0,v\Id_k)$ (for $s=0$, $T_0g$ is constant).

\emph{Gaussian calculus.} Let $C(h)$, $h\in(h_0,h_1)$, be a $C^1$ family of positive definite $K\times K$ matrices, and
let $H$ be a $C^2$ function on $\R^K$ which, with its derivatives up to order two, is bounded. Then
\begin{equation}\label{eq:interp}
\frac{d}{dh}\E_{C(h)}H(\zeta)=\frac12\sum_{a,b}\dot C_{ab}(h)\,\E_{C(h)}\big[\partial_a\partial_bH(\zeta)\big],
\end{equation}
by the heat equation $\partial_h\phi_{C(h)}=\frac12\sum_{a,b}\dot C_{ab}\partial_a\partial_b\phi_{C(h)}$ for the
Gaussian density (compare Fourier transforms) and two integrations by parts. If $\dot C$ has zero diagonal blocks and
$H(\zeta)=\prod_jG_j(\zeta_j)$ with $G_j\in\cS$, this reads
\begin{equation}\label{eq:drho}
\frac{d}{dh}\E_{C(h)}\Big[\prod_jG_j(\zeta_j)\Big]=\sum_{i<j}\E_{C(h)}\Big[\big\langle\nabla G_i(\zeta_i),
\dot C_{ij}\nabla G_j(\zeta_j)\big\rangle\prod_{l\ne i,j}G_l(\zeta_l)\Big].
\end{equation}

\emph{Entropy.} For a random variable $W>0$ with $\E W\log W<\infty$ put $\Ent(W)=\E W\log W-\E W\log\E W\ge0$. If
$\eta$ is a random vector, then $\Ent(W)=\E[\Ent(W\,|\,\eta)]+\Ent(\E[W\,|\,\eta])\ge\E[\Ent(W\,|\,\eta)]$. If
$W=\prod_jW_j$ with independent $W_j>0$, then $\Ent(W)=\sum_j\Ent(W_j)\prod_{l\ne j}\E W_l$ (expand $\log W$).

\section{A weighted reverse log-Sobolev inequality}
\begin{lemma}\label{lem:rlsi}
Let $\zeta$ be a centred Gaussian vector on $\R^K$ with covariance $C$, let $D=\mathrm{diag}(D_1,\dots,D_n)$ be
block-diagonal and positive semi-definite with $C\ge D$, and let $B$ be symmetric with zero diagonal blocks and
$D+B\ge0$. Then for $\Psi_j\in\cS$ on $\R^{k_j}$,
\begin{equation}\label{eq:rlsi}
\Ent\Big(\prod_j\Psi_j(\zeta_j)\Big)\ \ge\ -\sum_{i<j}\E\Big[\big\langle\nabla\Psi_i(\zeta_i),B_{ij}\nabla\Psi_j(\zeta_j)
\big\rangle\prod_{l\ne i,j}\Psi_l(\zeta_l)\Big].
\end{equation}
\end{lemma}

For $D=\lambda\Id$ this is~\cite[Lemma~1.6]{FR}; the proof is the same.

\begin{proof}
\emph{One Gaussian.} Let $\nu$ be the law of $m+A^{1/2}G$, $G\sim\gamma_k$, with $A\ge0$, and let $\Psi\in\cS$. For
$h(w)=\langle\theta,w-m\rangle$ the entropy inequality $\E_\nu[\Psi h]\le\Ent_\nu(\Psi)+\E_\nu\Psi\log\E_\nu e^h$
holds, $\E_\nu e^h=e^{\langle A\theta,\theta\rangle/2}$, and Gaussian integration by parts in $G$ gives
$\E_\nu[\Psi(w)(w-m)]=A\,\E_\nu\nabla\Psi$. With $\theta=\E_\nu\nabla\Psi/\E_\nu\Psi$ we obtain
\begin{equation}\label{eq:onegauss}
\Ent_\nu(\Psi)\ \ge\ \frac{\langle A\,\E_\nu\nabla\Psi,\E_\nu\nabla\Psi\rangle}{2\,\E_\nu\Psi}.
\end{equation}

\emph{Conditioning.} Write $\zeta=\eta+\omega$ with $\eta\sim N(0,C-D)$ and $\omega\sim N(0,D)$ independent. Since $D$
is block-diagonal, given $\eta$ the blocks $\zeta_j=\eta_j+\omega_j$ are independent with laws $N(\eta_j,D_j)$. Put
$u_j(\eta_j)=\E[\Psi_j(\zeta_j)\,|\,\eta]$, so that $u_j\ge\inf\Psi_j>0$, $u_j$ is $C^1$ and
$\nabla u_j(\eta_j)=\E[\nabla\Psi_j(\zeta_j)\,|\,\eta]$. Put $\beta_j=\nabla u_j/u_j$, $\beta=(\beta_1,\dots,\beta_n)$
and $U=\prod_ju_j(\eta_j)$. By the facts on entropy and~\eqref{eq:onegauss} with $A=D_j$,
\[
\Ent\Big(\prod_j\Psi_j(\zeta_j)\Big)\ \ge\ \E\Big[\sum_j\Ent\big(\Psi_j(\zeta_j)\,\big|\,\eta\big)\prod_{l\ne j}u_l(\eta_l)\Big]
\ \ge\ \frac12\,\E\big[U\,\beta^TD\beta\big].
\]
By conditional independence, and since $B$ has zero diagonal blocks, the sum on the right of~\eqref{eq:rlsi} without
the minus sign equals $\sum_{i<j}\E[U\langle\beta_i,B_{ij}\beta_j\rangle]=\frac12\E[U\beta^TB\beta]$. Finally
$\beta^TD\beta\ge-\beta^TB\beta$ pointwise, because $D+B\ge0$.
\end{proof}

\section{Proof of Theorem~\ref{thm:main}}
\subsection{The flow}
Let $f_j\in\cS$. For $q>0$ and $C\in\mathcal C$ put
\[
W=\prod_jf_j(\zeta_j)^q,\qquad\Phi(q,C)=\frac1q\log\E_CW,\qquad a_j=\frac{\nabla f_j}{f_j}(\zeta_j),\quad
a=(a_1,\dots,a_n).
\]
Let $C(h)$ be a $C^1$ family in $\mathcal C$ of positive definite matrices and $q(h)$ a $C^1$ positive function. Then
\begin{equation}\label{eq:deriv}
q^2\,\E W\,\frac{d}{dh}\Phi\big(q(h),C(h)\big)=\dot q\,\Ent(W)+\frac{q^3}2\,\E\big[W\,a^T\dot C a\big],
\end{equation}
all expectations with respect to $C(h)$. Indeed,
$\partial_q\Phi=-\frac1{q^2}\log\E W+\frac1q\E[W\sum_j\log f_j]/\E W=\Ent(W)/(q^2\E W)$, since
$\log W=q\sum_j\log f_j$. And $\dot C$ has zero diagonal blocks, so~\eqref{eq:drho} with $G_j=f_j^q$,
$\nabla G_j=qG_ja_j$, gives $\frac{d}{dh}\E_{C(h)}W=\frac{q^2}2\E[Wa^T\dot Ca]$ at fixed $q$. All functions involved
are bounded with bounded derivatives, so the differentiations are justified.

\subsection{The path}\label{sec:segment}
Let $D$ satisfy~\eqref{eq:lmi}. For $h\in[0,1]$ put
\[
C(h)=(1-h)\Sigma_0+h\Sigma_1,\qquad m(h)=1-h\Big(1-\frac1p\Big),\qquad q(h)=\frac1{m(h)} .
\]
Then $C(h)\in\mathcal C$, $C(h)\ge h\Sigma_1$ is positive definite for $h>0$, $q(0)=1$ and $q(1)=p$. Moreover:
\begin{enumerate}[label=(P\arabic*)]
\item $C(h)\ge m(h)D$. The difference is affine in $h$; at $h=0$ it is $\Sigma_0-D\ge0$, and at $h=1$ it is
$\Sigma_1-\frac1pD=\big(\Sigma_1-\Sigma_0+(1-\frac1p)D\big)+(\Sigma_0-D)\ge0$, by both halves of~\eqref{eq:lmi}.
\item With $\kappa(h)=\frac{p\,m(h)}{p-1}$ and $B=\kappa(h)(\Sigma_1-\Sigma_0)$, which has zero diagonal blocks,
$m(h)D+B=m(h)\big(D-\frac p{p-1}(\Sigma_0-\Sigma_1)\big)\ge0$.
\item $\dot q\,\kappa=q$: indeed $\dot q=(1-\frac1p)/m^2$, so $\dot q\kappa=\frac{(p-1)/p}{m^2}\cdot\frac{pm}{p-1}=\frac1m$.
\end{enumerate}

\subsection{Monotonicity}\label{sec:mono}
Let $h\in(0,1)$ and $\Lambda=q^2\E W\frac{d}{dh}\Phi(q(h),C(h))$. By~\eqref{eq:deriv} with $\dot C=\Sigma_1-\Sigma_0$,
\[
\Lambda=\dot q\,\Ent(W)+\frac{q^3}2\,G,\qquad G=\E\big[W\,a^T(\Sigma_1-\Sigma_0)a\big].
\]
Apply Lemma~\ref{lem:rlsi} with covariance $C(h)$, the block-diagonal matrix $m(h)D$ (admissible by (P1) and
Remark~\ref{rem:basic}(a)), the matrix $B$ of (P2), and $\Psi_j=f_j^q$. Since $\nabla\Psi_j=q\Psi_ja_j$, the right side
of~\eqref{eq:rlsi} is $-q^2\sum_{i<j}\E[W\langle a_i,B_{ij}a_j\rangle]=-\frac{q^2}2\E[Wa^TBa]=-\frac{\kappa q^2}2G$.
Hence $\Ent(W)\ge-\frac{\kappa q^2}2G$, whatever the sign of $G$, and since $\dot q>0$, (P3) gives
\[
\Lambda\ \ge\ \frac{q^2G}2\,(q-\dot q\kappa)=0 .
\]
So $h\mapsto\Phi(q(h),C(h))$ is non-decreasing on $(0,1)$. It is continuous on $[0,1]$: $q$ is continuous, the
integrands $\prod_jf_j^q$ are bounded and continuous, and $N(0,C(h))$ depends weakly continuously on $h$. Therefore
\[
\log\E_{\Sigma_0}\Big[\prod_jf_j(\xi_j)\Big]=\Phi(1,\Sigma_0)\le\Phi(p,\Sigma_1)=\log\Big\|\prod_jf_j(\xi_j)\Big\|_{L^p(\Sigma_1)},
\]
which is~\eqref{eq:comp} for $f_j\in\cS$.

\subsection{Approximation}\label{sec:approx}
Let $f_j:\R^{k_j}\to[\varepsilon,M]$ be measurable, $0<\varepsilon<M<\infty$, and $f_j^\delta=T_{1-\delta}f_j\in\cS$ for
$\delta\in(0,1)$. As $\delta\to0$, $f_j^\delta\to f_j$ in $L^2(\gamma_{k_j})$ ($T_s\to\Id$ strongly on $L^2$), hence in
$\gamma_{k_j}$-probability. Under both $\Sigma_0$ and $\Sigma_1$ each $\xi_j$ has law $\gamma_{k_j}$, so
$\prod_jf_j^\delta(\xi_j)\to\prod_jf_j(\xi_j)$ in probability, and all values lie in $[\varepsilon^n,M^n]$. Dominated
convergence carries~\eqref{eq:comp} from $f_j^\delta$ to $f_j$. For general measurable $f_j\ge0$ apply this to
$\max\{\varepsilon,\min\{f_j,M\}\}$, let $\varepsilon\downarrow0$ (dominated convergence, values in $[0,M+1]$), and
then $M\uparrow\infty$ (monotone convergence on both sides).

\subsection{Necessity}
For $f_j=e^{\langle a_j,\cdot\rangle}$, $\E_{\Sigma_0}e^{\langle a,\xi\rangle}=e^{a^T\Sigma_0a/2}$ and
$\|e^{\langle a,\xi\rangle}\|_{L^p(\Sigma_1)}=e^{pa^T\Sigma_1a/2}$. So~\eqref{eq:comp} forces $\Sigma_0\le p\Sigma_1$.
If $D$ satisfies~\eqref{eq:lmi}, then $D\le\Sigma_0$ and $\Sigma_1-\Sigma_0+(1-\frac1p)D\ge0$ give
$\Sigma_1-\frac1p\Sigma_0\ge\Sigma_1-\Sigma_0+(1-\frac1p)D\ge0$.\qed

\subsection{Curved paths}\label{sec:paths}
The argument uses neither the linearity of the path nor the constancy of $D$. We parametrise by $s=\log q$.

\begin{theorem}\label{thm:path}
Let $\Sigma_0,\Sigma_1\in\mathcal C$ and $T>0$. Let $C:[0,T]\to\mathcal C$ be continuous and piecewise $C^1$, with
finitely many pieces, $C(0)=\Sigma_0$, $C(T)=\Sigma_1$, and $C(s)$ positive definite for $0<s\le T$. Suppose that
for every $s$ in the interior of a piece there is a block-diagonal matrix $\tilde D(s)$ with
\begin{equation}\label{eq:pathcond}
0\le\tilde D(s)\le C(s),\qquad C'(s)+\tilde D(s)\ge0 .
\end{equation}
Then~\eqref{eq:comp} holds with $p=e^T$ for all measurable $f_j\ge0$. This remains true if, on some pieces, the
blocks are grouped: for a partition of $\{1,\dots,n\}$ into groups such that the diagonal blocks of $C(s)$ belonging to
the groups are constant on the piece, $\tilde D(s)$ may be block-diagonal for the groups.
\end{theorem}

\begin{proof}
Let $f_j\in\cS$ and $\tilde\Phi(s)=\Phi(e^s,C(s))$. In the interior of a piece,~\eqref{eq:deriv} with $q=e^s$ and
$\dot q=q$ gives $q^2\E W\,\tilde\Phi'(s)=q\Ent(W)+\frac{q^3}2\E[Wa^TC'(s)a]$. The matrix $C'(s)$ has zero diagonal
blocks, so Lemma~\ref{lem:rlsi} applies with covariance $C(s)$, the matrix $\tilde D(s)$, $B=C'(s)$ and
$\Psi_j=f_j^q$, and gives $\Ent(W)\ge-\frac{q^2}2\E[Wa^TC'(s)a]$. Hence $\tilde\Phi'(s)\ge0$. On a piece with grouped
blocks, $C'(s)$ has zero diagonal blocks for the groups, and Lemma~\ref{lem:rlsi}, which does not need identity
diagonal blocks, applies in the same way to the grouped blocks with $\Psi_G=\prod_{j\in G}f_j^q$. The function
$\tilde\Phi$ is continuous on $[0,T]$, also at a singular $\Sigma_0$: as $s\to s_0$, $\prod_jf_j^{e^s}\to\prod_jf_j^{e^{s_0}}$
uniformly, with values in a fixed compact subset of $(0,\infty)$, and $N(0,C(s))\to N(0,C(s_0))$ weakly. By the mean
value theorem on each closed piece, $\tilde\Phi(0)\le\tilde\Phi(T)$, which is~\eqref{eq:comp} with $p=e^T$ for
$f_j\in\cS$; Section~\ref{sec:approx} gives the general case.
\end{proof}

\begin{remark}\label{rem:path}
(a) Theorem~\ref{thm:main} is the case of the segment. With $q=e^s=1/m(h)$, the path of Section~\ref{sec:segment} has
$dC/ds=\kappa(\Sigma_1-\Sigma_0)$, and $\tilde D=m(h)D$ satisfies~\eqref{eq:pathcond} by (P1) and (P2).

(b) Along an admissible path, $C+C'=(C-\tilde D)+(C'+\tilde D)\ge0$, so $s\mapsto e^sC(s)$ is non-decreasing and
$\Sigma_0\le e^T\Sigma_1$. Let $T=\log p^*$ and $\Sigma_0v=p^*\Sigma_1v$. Then $e^sv^TC(s)v$ is non-decreasing with
equal values at $s=0$ and $s=T$, hence constant, so $v^T(C-\tilde D)v=0$ and $(C(s)-\tilde D(s))v=0$ in the interiors
of the pieces. For scalar blocks and $v$ without zero entries this forces $\tilde D(s)=\mathrm{diag}(C(s)v/v)$, and
letting $s\to0$ gives $\Sigma_0\ge\mathrm{diag}(\Sigma_0v/v)$. In Proposition~\ref{prop:example} this fails
($\Sigma_0-\mathrm{diag}(\frac54,\frac54,1)$ has the eigenvalue $-\frac12$), so no admissible path with scalar blocks
has $T=\log p^*$.

(c) For $n=2$ and $k_j=1$, write $c(s)$ for the correlation. Some $\tilde D(s)$ satisfies~\eqref{eq:pathcond} iff
$|c'(s)|\le1-|c(s)|$. Indeed $\tilde D=(1-|c|)\Id$ works then, and conversely
$|c'|^2\le\tilde D_{11}\tilde D_{22}=(1-u_1)(1-u_2)\le(1-\sqrt{u_1u_2})^2\le(1-|c|)^2$, where
$u_i=1-\tilde D_{ii}\in[0,1]$ and $u_1u_2\ge c^2$. So paths raise a correlation from $c_0\ge0$ to $c_1\in[c_0,1)$ iff
$p\ge\frac{1-c_0}{1-c_1}$, which is sharp (Corollary~\ref{cor:sharp}(a)), and lower it from $\theta\in(0,1)$ to $0$
iff $p\ge\frac1{1-\theta}$. The sharp exponent for lowering is $1+\theta$ (Remark~\ref{rem:forward}), while the
segment of Theorem~\ref{thm:main} needs $p\ge\frac{1-\theta}{1-2\theta}$ and $\theta<\frac12$.

(d) Along the segment $C(h)=(1-h)\Sigma_0+h\Sigma_1$ traversed at speed $dh/ds=1/\sigma(h)$, condition~\eqref{eq:pathcond}
reads $\sigma(h)\tilde D\ge\Sigma_0-\Sigma_1$. So~\eqref{eq:comp} holds with $p=\exp\int_0^1\sigma(h)\,dh$ whenever $\sigma>0$
is piecewise continuous and for each $h$ some block-diagonal $\tilde D$ with $0\le\tilde D\le C(h)$ satisfies
$\sigma(h)\tilde D\ge\Sigma_0-\Sigma_1$. Taking $\tilde D=m(h)D$ gives $\sigma=(1-\frac1p)/m(h)$, with integral $\log p$:
this is Theorem~\ref{thm:main}. For $n=2$ every admissible path is a reparametrised segment, so the gain in (c) over
Theorem~\ref{thm:main} comes entirely from letting $\tilde D$ vary.
\end{remark}

\begin{remark}\label{rem:groups}
In Proposition~\ref{prop:example} the grouped form of Theorem~\ref{thm:path} removes the obstruction of
Remark~\ref{rem:path}(b) at $s=0$: for the groups $G=\{1,2\}$ and $\{3\}$ one may take $\tilde D=\Sigma_0$. It does
not make $p^*$ reachable there. On a lossless path, $\tilde D(s)v=C(s)v$ and $(C'(s)+\tilde D(s))v=0$, so $C'v=-Cv$.
On a piece with the groups $G,\{3\}$ we have $C'_{GG}=0$, and comparing the two components of $C'v=-Cv$ gives
$v_G^TC_{GG}v_G=v_3^2$. On the first piece $C_{GG}=(\Sigma_0)_{GG}$, and
$v_G^T(\Sigma_0)_{GG}v_G=\frac52\ne\frac{21+\sqrt{41}}8=v_3^2$ for $v=(1,1,-\frac{1+\sqrt{41}}4)$. The other block
structures fail on the first piece as well: for the groups $\{1,3\},\{2\}$ the second component of
$\tilde D(s)v=C(s)v$ gives $\tilde D_{22}(s)=(C(s)v)_2/v_2\to\frac54$ as $s\to0$, while $\tilde D_{22}\le C_{22}=1$
(similarly for $\{2,3\},\{1\}$); and for a single group $C'=0$ on the piece, so the two identities give $C(s)v=0$,
which is impossible. Since the path has finitely many pieces, no admissible path, with any choice of block
structures, has $T=\log p^*$ in Proposition~\ref{prop:example}. The following proposition makes this quantitative.
\end{remark}

\begin{proposition}\label{prop:away}
In Proposition~\ref{prop:example} there is $\varepsilon>0$ such that every admissible path in the sense of
Theorem~\ref{thm:path}, with any block structures on its pieces, has $e^T\ge p^*+\varepsilon$.
\end{proposition}

\begin{proof}
For a partition $\sigma$ of $\{1,2,3\}$ into groups and $C\in\mathcal C$, let $\mathcal A_\sigma(C)$ be the set of pairs
$(V,\tilde D)$ with $V$ symmetric with zero diagonal blocks for $\sigma$, $\tilde D$ block-diagonal for $\sigma$,
$0\le\tilde D\le C$ and $V+\tilde D\ge0$. It is compact, since $\tilde D_{ii}\le C_{ii}=1$ and the $2\times2$ minors of
$V+\tilde D$ give $|V_{ij}|\le(\tilde D_{ii}\tilde D_{jj})^{1/2}\le1$, and $\{(C,V,\tilde D):(V,\tilde D)\in\mathcal A_\sigma(C)\}$
is closed. Let $v=(1,1,-\frac{1+\sqrt{41}}4)$, so $\Sigma_0v=p^*\Sigma_1v$, and put
\[
m(C)=\min\big\{v^T(C+V)v:\ \sigma,\ (V,\tilde D)\in\mathcal A_\sigma(C)\big\}.
\]
Since $v^T(C+V)v=v^T(C-\tilde D)v+v^T(V+\tilde D)v$, we have $m\ge0$; the minimum exists, since $(0,0)\in\mathcal A_\sigma(C)$, and $m$ is lower semicontinuous, since the
bounds above do not depend on $C$ and the graphs are closed.
Moreover $m(\Sigma_0)>0$: otherwise some $\sigma$ and $(V,\tilde D)\in\mathcal A_\sigma(\Sigma_0)$ have $(\Sigma_0-\tilde D)v=0$
and $(V+\tilde D)v=0$, i.e.\ $\tilde Dv=\Sigma_0v$ and $Vv=-\Sigma_0v$. For scalar blocks this forces
$\tilde D=\mathrm{diag}(\frac54,\frac54,1)\not\le\Sigma_0$; for the groups $\{1,2\},\{3\}$, the components of
$Vv=-\Sigma_0v$ give $v_G^T(\Sigma_0)_{GG}v_G=v_3^2$, which is false; for $\{1,3\},\{2\}$ (resp.\ $\{2,3\},\{1\}$) the second
(resp.\ first) component of $\tilde Dv=\Sigma_0v$ gives $\tilde D_{22}=\frac54$ (resp.\ $\tilde D_{11}=\frac54$),
contradicting $\tilde D\le\Sigma_0$; and for a single group $V=0$
gives $\Sigma_0v=0$ (all as in Remarks~\ref{rem:path}(b) and~\ref{rem:groups}). Choose $\delta\in(0,\log p^*]$ with
$m(C)\ge\eta:=\frac12m(\Sigma_0)$ whenever $C\in\mathcal C$ and $\max_{ij}|C_{ij}-(\Sigma_0)_{ij}|\le\delta$.

Along an admissible path, $(C'(s),\tilde D(s))\in\mathcal A_\sigma(C(s))$ in the interior of each piece, so
$|C'_{ij}|\le1$ and $\max_{ij}|C_{ij}(s)-(\Sigma_0)_{ij}|\le s$. The function $s\mapsto e^sv^TC(s)v$ is continuous,
piecewise $C^1$, with derivative $e^sv^T(C+C')v\ge0$; comparing its values at $0$ and $T$ gives $e^T\ge p^*$, so
$T\ge\delta$. Hence
\[
e^Tv^T\Sigma_1v-v^T\Sigma_0v=\int_0^Te^sv^T(C(s)+C'(s))v\,ds\ \ge\ \int_0^\delta m(C(s))\,ds\ \ge\ \eta\delta ,
\]
and $e^T\ge p^*+\eta\delta/(v^T\Sigma_1v)$. The argument does not give an explicit $\varepsilon$; numerically the best
value is about $0.016$ (see after Proposition~\ref{prop:certificate}).
\end{proof}

\section{Proofs of the corollaries}
\begin{lemma}\label{lem:forced}
Let $p=p^*>1$ and let $v\ne0$ satisfy $\Sigma_0v=p^*\Sigma_1v$. If $D$ satisfies~\eqref{eq:lmi} at $p=p^*$, then
$Dv=\Sigma_0v$.
\end{lemma}

\begin{proof}
Using $v^T\Sigma_1v=v^T\Sigma_0v/p^*$, we get $v^T\big(D-\frac{p^*}{p^*-1}(\Sigma_0-\Sigma_1)\big)v
=v^TDv-\frac{p^*}{p^*-1}\big(1-\frac1{p^*}\big)v^T\Sigma_0v=-v^T(\Sigma_0-D)v$. Both
$D-\frac{p^*}{p^*-1}(\Sigma_0-\Sigma_1)$ and $\Sigma_0-D$ are positive semi-definite, so $v^T(\Sigma_0-D)v=0$ and
hence $(\Sigma_0-D)v=0$.
\end{proof}

\begin{lemma}\label{lem:Z}
Let $A$ be a real symmetric $n\times n$ matrix with $A_{ij}\le0$ for $i\ne j$, and let $v\in\R^n$ have positive
entries with $Av=0$. Then $A\ge0$.
\end{lemma}

\begin{proof}
For $w\in\R^n$, using $A_{ii}v_i=-\sum_{j\ne i}A_{ij}v_j$ and the symmetry of $A$,
$w^TAw=-\frac12\sum_{i\ne j}A_{ij}v_iv_j\big(\frac{w_i}{v_i}-\frac{w_j}{v_j}\big)^2\ge0$.
\end{proof}

\begin{proof}[Proof of Corollary~\ref{cor:sharp}]
In each case necessity is part of Theorem~\ref{thm:main}, and for $p\ge p^*$ it suffices, by
Remark~\ref{rem:basic}(b), to find $D$ satisfying~\eqref{eq:lmi}, or $S$ satisfying~\eqref{eq:holder}, at
$p=p^*$, if $p^*>1$. If $p^*=1$ then $\Sigma_0=\Sigma_1$, and~\eqref{eq:comp} is Lyapunov's inequality.

(a) Let $\lambda=\lambda_{\min}(\Sigma_1)>0$. The eigenvalues of $\Sigma_0$ are $1-t+t\sigma$, $\sigma$ an eigenvalue
of $\Sigma_1$, with the same eigenvectors, and $(1-t+t\sigma)/\sigma$ decreases in $\sigma$. So $\Sigma_0\le p\Sigma_1$
iff $1-t+t\lambda\le p\lambda$. Assume this with $p>1$, and let $D=d\,\Id$ with $d=1-t+t\lambda$. Then
$\Sigma_0-D=t(\Sigma_1-\lambda\Id)\ge0$, and $\frac p{p-1}(\Sigma_0-\Sigma_1)=\frac p{p-1}(1-t)(\Id-\Sigma_1)
\le\frac p{p-1}(1-t)(1-\lambda)\Id$. The last bound is $\le d\,\Id$ iff
$p(1-t)(1-\lambda)\le(p-1)(1-t+t\lambda)$, i.e.\ iff $1-t+t\lambda\le p\lambda$.

(b) Conjugating by $\mathrm{diag}(\varepsilon_j)$ (that is, replacing $\xi_j$ by $\varepsilon_j\xi_j$ and $f_j$ by
$f_j(\varepsilon_j\,\cdot)$) we may assume $(\Sigma_1)_{ij}\le(\Sigma_0)_{ij}\le0$ for $i\ne j$. If
$(\Sigma_1)_{ij}=0$ then $(\Sigma_0)_{ij}=0$. So if the graph $\mathcal G$ on $\{1,\dots,n\}$ with edges
$\{(\Sigma_1)_{ij}\ne0\}$ is disconnected, the coordinates in different components are independent under both
measures, both sides of~\eqref{eq:comp} factorise over components, and $\Sigma_0\le p\Sigma_1$ holds iff it holds on
each component. So let $\mathcal G$ be connected, and let $p^*>1$. Put $N=p^*\Sigma_1-\Sigma_0\ge0$; it is singular.
For $i\ne j$, $N_{ij}=p^*(\Sigma_1)_{ij}-(\Sigma_0)_{ij}\le(\Sigma_1)_{ij}-(\Sigma_0)_{ij}\le0$, with
$N_{ij}<0$ whenever $(\Sigma_1)_{ij}<0$. So $sI-N$ is entrywise non-negative and irreducible for large $s$, and by the
Perron--Frobenius theorem its largest eigenvalue, $s-\lambda_{\min}(N)=s$, is simple with an eigenvector $v$ having
positive entries. Thus $Nv=0$, i.e.\ $\Sigma_0v=p^*\Sigma_1v$. Let $D=\mathrm{diag}\big((\Sigma_0v)_i/v_i\big)$.
Then $(\Sigma_0-D)v=0$, and $\Sigma_0-D$ has non-positive off-diagonal entries, so $\Sigma_0-D\ge0$ by
Lemma~\ref{lem:Z}. Next, $A'=\Sigma_1-\Sigma_0+(1-\frac1{p^*})D$ has the off-diagonal entries
$(\Sigma_1-\Sigma_0)_{ij}\le0$, and $A'v=\Sigma_1v-\frac1{p^*}\Sigma_0v=0$; so $A'\ge0$ by Lemma~\ref{lem:Z}.
Multiplying by $\frac{p^*}{p^*-1}$ gives the left half of~\eqref{eq:lmi} at $p=p^*$.

(c) The eigenvalues of a symmetric matrix $\begin{psmallmatrix}a\Id&B\\B^T&a\Id\end{psmallmatrix}$ are among $a$ and
$a\pm\sigma_i$, $\sigma_i$ the singular values of $B$, and $a\pm\sigma_{\max}$ occur. Hence $\lambda_{\max}(\Sigma_0)=1+|t|\,\|C\|$,
$\lambda_{\min}(\Sigma_1)=1-\|C\|$, and $p\Sigma_1-\Sigma_0\ge0$ iff $|p-t|\,\|C\|\le p-1$. For $t\le0$ this reads
$1+|t|\,\|C\|\le p(1-\|C\|)$, i.e.\ $\lambda_{\max}(\Sigma_0)\le p\lambda_{\min}(\Sigma_1)$, and
Proposition~\ref{prop:holder} applies.
\end{proof}

\begin{proof}[Proof of Corollary~\ref{cor:IK}]
Let $A_j\sim\gamma_{k_j}$ be independent of each other and of $\xi\sim N(0,\Sigma)$. Then
$\prod_jT_{x_j}f_j(\xi_j)=\E[\prod_jf_j(x_j\xi_j+\sqrt{1-x_j^2}A_j)\,|\,\xi]$, and the vector
$(x_j\xi_j+\sqrt{1-x_j^2}A_j)_j$ has covariance $X\Sigma X+\Id-X^2=\Sigma_0\in\mathcal C$. By Tonelli,
$\E_\Sigma[\prod_jT_{x_j}f_j(\xi_j)]=\E_{\Sigma_0}[\prod_jf_j(\xi_j)]$. With $p_j=p$ and $\alpha=1/p$,
\eqref{eq:IK} is the $p$-th power of~\eqref{eq:comp}, and~\eqref{eq:IKcond} reads
$\frac1p(\Id-X^2)\le\Sigma-\frac1pX\Sigma X$, i.e.\ $\Sigma_0\le p\Sigma$. Positive measurable $f_j$ are covered by
Theorem~\ref{thm:main} and Proposition~\ref{prop:holder}, and the exponential witnesses of necessity are positive.
The off-diagonal blocks of $\Sigma_0$ are $(\Sigma_0)_{ij}=x_ix_j\Sigma_{ij}$. In case (i) this gives
$\Sigma_0=(1-t)\Id+t\Sigma$ with $t\in[0,1]$, and Corollary~\ref{cor:sharp}(a) applies. In case (ii),
$\varepsilon_i\varepsilon_j\Sigma_{ij}\le\varepsilon_i\varepsilon_jx_ix_j\Sigma_{ij}\le0$ because $x_ix_j\in[0,1]$ whenever
$\Sigma_{ij}\ne0$; this is~\eqref{eq:signs}, and Corollary~\ref{cor:sharp}(b) applies. In case (iii), $\Sigma_0$ has
the off-diagonal block $t\Sigma_{12}$ with $t=x_1x_2$, and Corollary~\ref{cor:sharp}(a) (if $t\ge0$) or (c) (if
$t<0$) applies. In case (iv), $(\Sigma_0)_{ij}=\Sigma_{ij}$ for $i,j\in G$. If $x_j=0$ for all $j\notin G$, take the
groups $G$ and $\{j\}$, $j\notin G$: all blocks of $\Sigma_0$ between different groups vanish, which is the grouped
form of Corollary~\ref{cor:sharp}(a) with $t=0$. If $G$ contains all indices but $j_0$, take the two groups $G$ and
$\{j_0\}$: the block of $\Sigma_0$ between them is $x_{j_0}$ times that of $\Sigma$, which is the grouped form of
(a) or (c) with $t=x_{j_0}$.
\end{proof}

\begin{proof}[Proof of Proposition~\ref{prop:example}]
Exact computation (checked with sympy, \texttt{numerics/comparison\_example.py}): $\det(\Sigma_0-p\Sigma_1)
=-\frac1{16}(2p-3)(4p^2-11p+5)$, so the generalised eigenvalues are $\frac32$ and $\frac{11\pm\sqrt{41}}8$, and
$p^*=\frac{11+\sqrt{41}}8$. Both matrices are invariant under swapping the first two coordinates, and the sets of
solutions of~\eqref{eq:lmi} and~\eqref{eq:holder} are convex. Averaging a solution with its swap, we may take
$D=\mathrm{diag}(d,d,e)$ and $S=\mathrm{diag}(s,s,s_3)$.

Condition~\eqref{eq:lmi}: $\Sigma_0-D\ge0$ iff $d\le\frac34$ and $e\le1$ (eigenvalues $1-d\pm\frac14$ and $1-e$).
With $c=\frac p{p-1}$, the matrix $D-c(\Sigma_0-\Sigma_1)$ has the eigenvector $(1,-1,0)$ with eigenvalue
$d-\frac c4$, and on the span of $\frac1{\sqrt2}(1,1,0)$ and $(0,0,1)$ it acts as
$\begin{psmallmatrix}d+c/4&c/\sqrt2\\c/\sqrt2&e\end{psmallmatrix}$. So it is positive semi-definite iff
$d\ge\frac c4$ and $(d+\frac c4)e\ge\frac{c^2}2$. These constraints become weaker as $d$ and $e$ grow, so~\eqref{eq:lmi}
is feasible iff they hold at $d=\frac34$, $e=1$, i.e.\ iff $c\le3$ and $3+c\ge2c^2$. Since $(2c-3)(c+1)\le0$ means
$c\le\frac32$, this is $p\ge3$. At $p=3$, $D=\mathrm{diag}(\frac34,\frac34,1)$: $\Sigma_0-D$ has the eigenvalues
$\frac12,0,0$, and $D-\frac32(\Sigma_0-\Sigma_1)$ has the eigenvalues $\frac{17}8,\frac38,0$.

Condition~\eqref{eq:holder}: $S-\Sigma_0\ge0$ iff $s\ge\frac54$ and $s_3\ge1$, and $p\Sigma_1-S$ has the eigenvector
$(1,-1,0)$ with eigenvalue $\frac p2-s$. So~\eqref{eq:holder} forces $\frac54\le s\le\frac p2$, hence $p\ge\frac52$.
For $p\ge\frac52$ take $S=\mathrm{diag}(\frac54,\frac54,1)$: then $S-\Sigma_0$ has the eigenvalues $\frac12,0,0$, and at
$p=\frac52$ the matrix $p\Sigma_1-S$ has the eigenvalues $0$ and $2\pm\frac34\sqrt6$, all $\ge0$; for larger $p$ add
$(p-\frac52)\Sigma_1\ge0$.
\end{proof}

\begin{remark}\label{rem:forward}
For $\Sigma_1=\Id$, \eqref{eq:comp} is the forward Brascamp--Lieb (H\"older) inequality
$\E_{\Sigma_0}\prod_jf_j\le\prod_j\|f_j\|_{L^p(\gamma_{k_j})}$, which holds iff $\lambda_{\max}(\Sigma_0)\le p$
(\cite{CDP15}, see also~\cite{Ledoux14}; this is Proposition~\ref{prop:holder} with $S=p\Id$). Condition~\eqref{eq:lmi}
is not sharp there. Let $n=2$, $k_j=1$ and $\mathrm{cov}_{\Sigma_0}(\xi_1,\xi_2)=\theta\in(0,1)$, so that $p^*=1+\theta$.
By symmetry and convexity we may take $D=d\Id$, and~\eqref{eq:lmi} reads $c\theta\le d\le1-\theta$ with
$c=\frac p{p-1}$. So~\eqref{eq:lmi} is feasible iff $\theta<\frac12$ and $p\ge\frac{1-\theta}{1-2\theta}$, a
threshold larger than $1+\theta$; for $\theta\ge\frac12$ it is never feasible. Thus~\eqref{eq:lmi} is not necessary
for~\eqref{eq:comp}. The reason is that Lemma~\ref{lem:rlsi} is a reverse inequality.
\end{remark}

\section{Numerical checks}\label{sec:num}
All of the following is floating point, evidence for the algebra and not part of the proofs.
\texttt{numerics/check\_comparison.py} ($n=3$, $k_j=1$, tensor Gauss--Hermite quadrature) checks (A) Lemma~\ref{lem:rlsi}
for $300$ random covariances, diagonal $D$ and pairings $B$, with random smooth positive functions (smallest
ratio of the two sides about $1.8$); (B) the monotonicity along the path for $60$ random pairs $(\Sigma_0,\Sigma_1)$
of the form of Corollary~\ref{cor:IK} at which~\eqref{eq:lmi} is feasible (no decrease); and (C) the inequality at
$p=p^*$ for $25$ such pairs with optimised test functions: the optimisers found with $28$ nodes per axis exploit the
quadrature error, and recomputed with $60$ nodes per axis they give log-ratios $\le0$.
\texttt{numerics/comparison\_example.py} checks Proposition~\ref{prop:example} exactly (including both thresholds), confirms
the threshold $p=3$ for~\eqref{eq:lmi} by bisection, finds no violation of~\eqref{eq:comp} at $p=p^*$
by Gaussian test functions $e^{a_jy-b_jy^2/2}$ (closed forms), and prints the smallest eigenvalues of
$p^*\Sigma_1^{\circ d}-\Sigma_0^{\circ d}$ of Remark~\ref{rem:second} ($0$ for $d=1$, and $0.6102$, $0.8982$, $1.0381$ for
$d=2,3,4$). \texttt{numerics/comparison\_example\_attack.py} optimises, in the same example, over test
functions $f_j=\exp(a_jt+b_jt^2+c_j\tanh(w_j(t-d_j)))$ with bounded coefficients at $p\in\{p^*,2.3,2.6,2.9\}$ (only
$p^*$ lies in the open range): the best log-ratios are at most $2\cdot10^{-8}$ with $30$ quadrature nodes per axis,
and negative when recomputed with $56$. \texttt{numerics/comparison\_example\_attack2.py} uses
$f_j=\exp(\sum_{k\le6}a_{jk}T_k(s))$, $T_k$ Chebyshev polynomials, $s=t/\sqrt{16+t^2}$, $|a_{jk}|\le1.5$, optimised with
$36$ nodes per axis and re-evaluated with $56$ and $72$: at $p=p^*$ and $p=2.2$ all $30$ optimisers are stable, with best
log-ratios $-1.5\cdot10^{-7}$ and $-1.4\cdot10^{-7}$; a variant with Hermite polynomials and $|a_{jk}|\le0.3$ gives
$-1.8\cdot10^{-8}$ and $-2.2\cdot10^{-12}$ (\texttt{comparison\_example\_attack2\_he.py}). \texttt{numerics/comparison\_path\_flow.py} solves the discretised path
problem of Theorem~\ref{thm:path} (piecewise-linear paths, one diagonal $\tilde D$ per piece; a semidefinite
programme for fixed $T$, bisection in $T$). It reproduces the exponents of Remark~\ref{rem:path}(c) up to the
discretisation error, and in Proposition~\ref{prop:example} it gives $2.4405$, $2.3596$, $2.3224$, $2.3045$, $2.2957$
for $4,8,16,32,64$ pieces (and $2.2913$, $2.2891$ for $128$, $256$).
\texttt{numerics/comparison\_path\_twostage.py} and \texttt{comparison\_path\_multistage.py} do the same with the
groups $\{1,2\},\{3\}$ on some stretches (Theorem~\ref{thm:path}): two discretised stretches give $2.2152$ with $8+24$
pieces, and with $32+32$ pieces the problem is feasible at $p=2.21$ but not at $2.20$; alternating three or four
stretches ($16$ and $12$ pieces each, probed at $p=2.20$ and $2.21$) was not better at the splits we tried. \texttt{numerics/comparison\_chain\_scan.py} scans the construction
of Proposition~\ref{prop:certificate} over $b$. Only the exact verifications
in Proposition~\ref{prop:certificate} are part of a proof. For random $n=3$ and $n=4$ instances of Corollary~\ref{cor:IK}
where~\eqref{eq:lmi} fails at $p^*$, Gaussian test functions do not violate~\eqref{eq:comp} at $p^*$ either ($60$ cases
each), and the matrix $D=\mathrm{diag}(\Sigma_0v/v)$ forced by Lemma~\ref{lem:forced} ($v$ had no zero entries in any
instance) satisfies~\eqref{eq:lmi} in all $2000$ random instances with sign pattern~\eqref{eq:signs} for each
$n\in\{3,4,6\}$, and in none of the random equicorrelated instances with $\rho>0$, $n\ge3$ and unequal parameters
(\texttt{numerics/comparison\_gaussian\_tests.py}).

\section{Novelty and limitations}
\emph{Novelty.} Theorem~\ref{thm:main} compares a Gaussian expectation under one covariance with a correlated $L^p$
norm under another, for product functions. Its special cases include the reverse H\"older inequality
of~\cite{CDP15} with equal exponents and Lyapunov's inequality, and Proposition~\ref{prop:holder} is an immediate
consequence of~\cite{CDP15}. Applied to the forward inequality of~\cite{IK25} with $\alpha=1/p$, a case excluded
from~\cite[Thm.~9(i)]{IK25}, the results recover the equal-parameter case~\cite[Thm.~1.1]{FR} at $r=1$ and, for two
$\rho$-correlated functions, \cite[Prop.~6.1]{FR}. What is new is Theorem~\ref{thm:main} itself and its path
version Theorem~\ref{thm:path}, Corollary~\ref{cor:sharp}(b), case (ii) of Corollary~\ref{cor:IK} with $n\ge3$ and unequal
parameters, and Propositions~\ref{prop:certificate}, \ref{prop:notsuff}, \ref{prop:gauss} and~\ref{prop:away}.
Gaussian saturation is known for many inequalities of Brascamp--Lieb type~\cite{Lieb90,Bar98,CDP15,CL21,BW22,NT25}.
Comparisons of $\E_{\Sigma_0}F$ with $\E_{\Sigma_1}F$ for a single function (the case $p=1$) under sign conditions on
mixed second derivatives go back to Slepian~\cite{Sle62} and Kahane~\cite{Kah86}, with the interpolation
formula~\eqref{eq:interp}. Milman, Nakamura and Tsuji~\cite{MNT26} prove inequalities for convex sets under Gaussian
measures with different covariance matrices, using a Gaussian forward--reverse Brascamp--Lieb inequality for
log-concave functions; these results concern convex sets and log-concave functions and do not
contain~\eqref{eq:comp} for general non-negative $f_j$. We did not find a statement covering~\eqref{eq:comp} with two
different covariances, but we did not search the literature systematically beyond these works
and~\cite{Ledoux14}. Novelty is claimed to the best of this search. The flow generalises the correlation flow
of~\cite[Section~4]{FR} in two ways: the covariance moves along an arbitrary segment rather than a ray through $\Id$,
and Lemma~\ref{lem:rlsi} conditions on a block-diagonal rather than a scalar Gaussian component.

\emph{Limitations.} (1) Condition~\eqref{eq:lmi} is sufficient but not necessary for~\eqref{eq:comp}: in
Remark~\ref{rem:forward},~\eqref{eq:comp} holds at $p^*=1+\theta$ while~\eqref{eq:lmi} fails there. The necessary
condition $\Sigma_0\le p\Sigma_1$ is not sufficient either (Proposition~\ref{prop:notsuff}), and we do not know the
exact threshold in general. The flows of this note (including the
construction of Proposition~\ref{prop:certificate}) do not reach the necessary condition in
Proposition~\ref{prop:example}: they give $p\ge2.1925$
(Proposition~\ref{prop:certificate}) against $p^*=2.175\ldots$, the range $[p^*,2.1925)$ is open, and by
Proposition~\ref{prop:away} the flows are bounded away from $p^*$ there. In the example $\theta=0.9$, $\rho=0.5$ after
Proposition~\ref{prop:notsuff} they give only $p\ge5$ (Proposition~\ref{prop:holder} gives $p\ge3.8$), while a
violation is found at $p=1.56$ and none for $p\ge1.6$ (floating point). For the forward
inequality~\eqref{eq:IK} with $\alpha=1/p$ this leaves open, for example, three positively equicorrelated scalar
functions with parameters not covered by Corollary~\ref{cor:IK}(i) and (iv), such as $x=(1,\frac12,0)$. For scalar
blocks the flow is rigid at $p=p^*$: if some $v$ with $\Sigma_0v=p^*\Sigma_1v$ has no zero entries, then
Lemma~\ref{lem:forced} leaves only the candidate $D=\mathrm{diag}(\Sigma_0v/v)$. (2) Theorems~\ref{thm:main} and~\ref{thm:path} are
not sharp in forward situations (Remarks~\ref{rem:path}(c) and~\ref{rem:forward}), and Proposition~\ref{prop:holder}
is not sharp on the segment of Corollary~\ref{cor:sharp}(a); we have no method that interpolates between the two. (3) Only the outer exponent
$\alpha=1/p$ of~\eqref{eq:IK} is covered. For other $\alpha\in(0,1)$ and two functions with unequal parameters,
condition~\eqref{eq:IKcond} is not sufficient in general when $r=\alpha p>2$~\cite[Props.~6.3, 6.5]{FR}. For $r>1$
and $\rho,x_1,x_2>0$ a flow with separate noise parameters gives a sufficient condition under which it is sharp,
verified at explicit points with $1<r\le2$~\cite[Thm.~6.9, Prop.~6.11]{FR}; it is also sufficient when $r\ge1$ and
$x_1x_2=0$~\cite[Prop.~6.4]{FR}. For $r<1$ the analogous flows cannot be sharp~\cite[Remark~6.7(b)]{FR}.

\begin{thebibliography}{99}
\bibitem{Bar98} F.~Barthe, On a reverse form of the Brascamp--Lieb inequality, Invent.\ Math.\ 134 (1998), 335--361.
\bibitem{BW22} F.~Barthe, P.~Wolff, Positive Gaussian kernels also have Gaussian minimizers, Mem.\ Amer.\ Math.\
Soc.\ 276 (2022), no.~1359; arXiv:1805.02455.
\bibitem{CDP15} W.-K.~Chen, N.~Dafnis, G.~Paouris, Improved H\"older and reverse H\"older inequalities for
Gaussian random vectors, Adv.\ Math.\ 280 (2015), 643--689.
\bibitem{CL21} T.~A.~Courtade, J.~Liu, Euclidean forward--reverse Brascamp--Lieb inequalities: finiteness, structure,
and extremals, J.\ Geom.\ Anal.\ 31 (2021), 3300--3350.
\bibitem{FR} Companion note: \emph{Correlated forward hypercontractivity with outer exponent $\alpha<1$}
(\texttt{RESULT\_forward\_real\_hypercontractivity.pdf}, this repository).
\bibitem{IK25} P.~Ivanisvili, P.~Kalantzopoulos, Multiversion of the Hausdorff--Young inequality,
arXiv:2506.08494 (2025).
\bibitem{Kah86} J.-P.~Kahane, Une in\'egalit\'e du type de Slepian et Gordon sur les processus gaussiens, Israel J.\
Math.\ 55 (1986), 109--110.
\bibitem{Ledoux14} M.~Ledoux, Remarks on Gaussian noise stability, Brascamp--Lieb and Slepian inequalities, in
Geometric Aspects of Functional Analysis, Lecture Notes in Math.\ 2116, Springer, 2014, 309--333.
\bibitem{Lieb90} E.~H.~Lieb, Gaussian kernels have only Gaussian maximizers, Invent.\ Math.\ 102 (1990),
179--208.
\bibitem{MNT26} E.~Milman, S.~Nakamura, H.~Tsuji, Gaussian inequalities with different covariance matrices -- the
heterovariant Gaussian Brunn--Minkowski inequality, arXiv:2610.13596 (2026).
\bibitem{NT25} S.~Nakamura, H.~Tsuji, A generalized Legendre duality relation and Gaussian saturation, Invent.\
Math.\ (2025), doi:10.1007/s00222-025-01382-5; arXiv:2409.13611.
\bibitem{Sle62} D.~Slepian, The one-sided barrier problem for Gaussian noise, Bell System Tech.\ J.\ 41 (1962),
463--501.
\end{thebibliography}
\end{document}
